% V40 validation-only manuscript draft. Compile from this directory.
% Evidence: V40-v2 expanded-adjacency component-disjoint validation only.
\documentclass[pdflatex,sn-basic,Numbered,iicol]{../tex/sn-jnl}

\usepackage{graphicx}
\usepackage{xcolor}
\usepackage{booktabs}
\usepackage{tabularx}
\usepackage{array}
\usepackage{multirow}
\usepackage{amsmath,amssymb,amsfonts}
\usepackage{enumitem}
\usepackage{placeins}
\hypersetup{hidelinks}

\newcolumntype{Y}{>{\raggedright\arraybackslash}X}
\setlist[itemize]{leftmargin=*,nosep,topsep=1pt}
\newcommand{\figref}[1]{Fig.~\ref{#1}}
\newcommand{\tabref}[1]{Table~\ref{#1}}

\title[Reliability-aware RGB-thermal-event fusion]{Reliability-Aware RGB-Thermal-Event Fusion for Lightweight UAV Vehicle Detection with Expanded-Adjacency Component-Disjoint Validation}

\author*[1]{\fnm{Nan} \sur{Xin}}\email{floomeer83felix@gmail.com}
\author[1]{\fnm{Xueting} \sur{Jin}}
\affil*[1]{\orgname{Anhui Police College}, \orgaddress{\city{Hefei}, \state{Anhui}, \country{China}}}

\abstract{RGB, thermal, and event observations can provide complementary information for UAV vehicle detection, but apparent gains can be overstated when visually adjacent observations cross a training-validation boundary. We present a lightweight reliability-aware RGB-thermal-event detector built from modality-specific stems, a softmax fusion gate, and a RepViT-FPN-FCOS detector. The reliability-aware configuration uses modality dropout with a pre-specified probability of 0.15; it is compared with a matched early-fusion baseline using the same downstream detector, training recipe, and two fixed random seeds. We construct a V40-v2 validation protocol by assigning connected components of a fixed graph to one partition, where graph edges arise from exact RGB matches, perceptual-hash candidates, and human-adjudicated adjacent-or-near-identical observations. The resulting validation partition contains 2,213 images and 5,867 vehicle boxes. Reliability-aware fusion improves two-run mean AP50 from 0.9408 to 0.9586, AP75 from 0.8208 to 0.8760, and F1 from 0.8945 to 0.9133. Image-level bootstrap intervals for the reliability-minus-early differences remain positive for AP50, AP75, and F1. Under deterministic synthetic removal of RGB, thermal, or event channels, the reliability-aware system retains higher AP50 than matched early fusion. The evidence is limited to one local tri-modal representation, a validation-only protocol, two fixed checkpoints per model group, and synthetic channel removal; it does not establish independent-test performance, external generalization, or physical sensor-failure robustness.}

\keywords{UAV vehicle detection; RGB-thermal-event fusion; reliability-aware fusion; modality dropout; component-disjoint validation; multi-sensor perception.}

\begin{document}
\raggedbottom
\maketitle

\section{Introduction}
UAV vehicle detection must localize small targets under changing altitude, viewpoint, scale, illumination, motion, and clutter. These conditions are challenging for generic object detection and for aerial imaging in particular \cite{zhao2019object,cheng2016survey,xia2022dota,zhu2017deepRS}. RGB imagery provides texture and color, but may weaken under poor contrast or changing exposure. Thermal observations can provide temperature-related contrast, while event representations can encode brightness changes with high temporal resolution and wide dynamic range \cite{baltrusaitis2019multimodal,gallego2022event}. The key design question is therefore not merely how to concatenate modalities, but how to use their complementary information without unnecessarily increasing the cost of a lightweight detector.

A second issue is evaluation integrity. In multimodal UAV recordings, neighboring observations may be visually similar even when their file contents are not identical. If such observations fall on different sides of a split, validation performance can be optimistic. This risk is distinct from model design and should be handled before interpreting a fusion gain \cite{kapoor2023leakage,cawley2010overfitting}. In the present project, an earlier split review identified exact RGB duplicates and subsequently identified cross-partition adjacent-or-near-identical observations through a human review workflow. These findings motivate a more conservative validation protocol rather than a stronger claim about universal independence.

This paper presents RA-RepDet, a compact five-channel RGB-thermal-event detector. The reliability-aware model uses separate stems for RGB, thermal, and event inputs, predicts image-level softmax fusion weights, and passes the fused representation through a RepViT-M0.9 \cite{wang2024repvit}, feature-pyramid network (FPN) \cite{lin2017fpn}, and FCOS detector \cite{tian2019fcos}. The matched early-fusion baseline differs only in its front-end fusion mechanism. During reliability-aware training, independent modality dropout reduces reliance on a single input group; the $p=0.15$ setting was specified from archived development evidence before V40 results were viewed. V40 does not perform a dropout sweep or select an optimal dropout rate.

The contribution of the work is twofold. First, it offers a lightweight, reliability-aware fusion front end for a shared detector stack. Second, it reports the result on a validation partition constructed to keep all components of a fixed candidate-and-adjacency graph within one side of the split. The conclusions are intentionally narrow: they concern a fixed validation protocol and controlled synthetic channel removal, not independent testing or real sensor-failure deployment.

\noindent\textbf{Contributions.}
\begin{itemize}
    \item We develop a lightweight tri-modal fusion front end with modality-specific stems, image-level softmax fusion weights, and a shared RepViT-FPN-FCOS detector.
    \item We compare matched early fusion and a pre-specified reliability-aware $p=0.15$ configuration under a frozen V40-v2 expanded-adjacency component-disjoint validation protocol.
    \item We report two-seed aggregate results, image-level bootstrap intervals, deterministic synthetic channel-removal results, and a reproducible efficiency measurement with explicit latency boundaries.
    \item We document the scope and limitations of the validation evidence, including the absence of an independent test set, trained single-modality baselines, and a V40 dropout sweep.
\end{itemize}

\section{Related Work}
\subsection{Aerial Detection and Lightweight Detectors}
Aerial detection involves wide scale variation, small objects, geometric changes, and cluttered backgrounds \cite{cheng2016survey,zhu2017deepRS,xia2022dota}. Modern detectors include region-proposal and dense one-stage formulations \cite{ren2017faster,lin2020focal}. FCOS is an anchor-free dense detector that predicts object locations directly \cite{tian2019fcos}. In this study, the detector head is held fixed; the research question is whether the tri-modal fusion front end can improve a compact detection pipeline. RepViT is used as a mobile-oriented backbone, and FPN provides multi-scale features \cite{wang2024repvit,lin2017fpn}.

\subsection{Multimodal and Event-Based Perception}
Multimodal learning integrates complementary sources while accounting for differences in their reliability and representation \cite{baltrusaitis2019multimodal,ramachandram2017deep}. Visible-infrared fusion has been studied extensively at pixel and feature levels \cite{li2017pixel,ma2019surveyfusion,li2019densefuse}. Event cameras provide asynchronous brightness-change information and can support perception under fast motion or high dynamic range \cite{lichtsteiner2008sensor,gallego2022event,gehrig2021dsec}. RA-RepDet does not reconstruct raw event streams or infer a missing modality. It consumes a preconstructed five-channel sample and learns a compact image-level mixture of modality-stem features.

\subsection{Missing Modality and Evaluation Validity}
Modality dropout is a common strategy for limiting brittle dependence on one information source \cite{neverova2016moddrop}. Here, it is used as a training-time intervention, while zeroing one channel group at evaluation time gives a controlled synthetic channel-removal protocol. This protocol is useful for comparing models under the same intervention, but it cannot reproduce sensor drift, misregistration, field-of-view change, or other physical failure modes.

Evaluation leakage and adaptive reuse of validation information can produce optimistic performance estimates \cite{kapoor2023leakage,cawley2010overfitting}. The present work does not claim that any finite audit proves the absence of all dependence. Instead, it constructs a split whose train-validation boundary does not cut components of a specified graph formed from exact RGB, pHash, dHash, and human-adjudicated adjacent-or-near-identical relations.

\section{Method}
\subsection{Problem Formulation and Detector Interface}
Each observation is a five-channel tensor
\begin{equation}
\mathbf{x}=[\mathbf{x}_{\mathrm{rgb}},\mathbf{x}_{\mathrm{th}},\mathbf{x}_{\mathrm{ev}}],
\end{equation}
where $\mathbf{x}_{\mathrm{rgb}}\in\mathbb{R}^{3\times H\times W}$, $\mathbf{x}_{\mathrm{th}}\in\mathbb{R}^{1\times H\times W}$, and $\mathbf{x}_{\mathrm{ev}}\in\mathbb{R}^{1\times H\times W}$. The task is single-class vehicle detection. Raw TriAir class 0 is mapped to the foreground label used by the detector, with background retained as label 0.

The downstream detector is common to both systems: a RepViT-M0.9 backbone, an FPN with 128 output channels, and an FCOS head. The common detector interface makes the comparison a system-level contrast between two fusion/training choices rather than a comparison of unrelated architectures.

\begin{figure*}[t]
\centering
\includegraphics[width=0.98\textwidth]{figures/fig1_method.pdf}
\caption{RA-RepDet. RGB, thermal, and event inputs are first encoded by modality-specific stems. A small image-level reliability module produces softmax fusion weights, the weighted feature is projected to the shared detector interface, and a RepViT-FPN-FCOS stack produces vehicle detections. Modality dropout is used only during training.}
\label{fig:method}
\end{figure*}

\subsection{Matched Early Fusion}
The matched early-fusion baseline first projects the five input channels to three channels using a learned $1\times1$ convolution. The resulting representation enters the same RepViT-FPN-FCOS detector. It therefore has access to all five channels but does not form explicit modality-specific features or sample-dependent modality weights.

\subsection{Reliability-Aware Fusion}
The reliability-aware model processes RGB, thermal, and event inputs through separate $3\times3$ convolution--batch-normalization--SiLU stems. Each stem produces a 16-channel feature map, denoted by $\mathbf{F}_{\mathrm{rgb}}$, $\mathbf{F}_{\mathrm{th}}$, and $\mathbf{F}_{\mathrm{ev}}$. Global average pooling yields one descriptor per stem. The descriptors are concatenated and passed through a small two-layer multilayer perceptron to generate three logits. The fusion weights are
\begin{equation}
\boldsymbol{\alpha}=\operatorname{softmax}\left(g\left([\operatorname{GAP}(\mathbf{F}_{\mathrm{rgb}});\operatorname{GAP}(\mathbf{F}_{\mathrm{th}});\operatorname{GAP}(\mathbf{F}_{\mathrm{ev}})]\right)\right),
\end{equation}
where $\boldsymbol{\alpha}=[\alpha_{\mathrm{rgb}},\alpha_{\mathrm{th}},\alpha_{\mathrm{ev}}]$ and the three entries sum to one. The fused feature is
\begin{equation}
\mathbf{F}_{\mathrm{fuse}}=\alpha_{\mathrm{rgb}}\mathbf{F}_{\mathrm{rgb}}+\alpha_{\mathrm{th}}\mathbf{F}_{\mathrm{th}}+\alpha_{\mathrm{ev}}\mathbf{F}_{\mathrm{ev}}.
\end{equation}
A projection converts $\mathbf{F}_{\mathrm{fuse}}$ to the three-channel interface expected by the shared backbone. The gate is a learned feature-mixture mechanism; its weights are not calibrated probabilities of physical sensor health.

\subsection{Modality-Dropout Training}
For the reliability-aware model, RGB channels $0{:}3$, thermal channel 3, and event channel 4 are independently zeroed during training with probability $p=0.15$. If all three groups are selected, one group is restored uniformly at random so that at least one stream remains. The $p=0.15$ setting was fixed for the V40 comparison before the V40 result was examined. Ordinary full-modality inference supplies all available streams.

\section{V40-v2 Dataset and Validation Protocol}
The local TriAir representation contains 10,489 five-channel samples and 30,634 valid vehicle boxes. It retains 738 observations without label text files and one empty label text file as empty-target observations. Public URL, license, version, redistribution terms, and independently verified temporal metadata were not established in this study.

The V40-v2 split begins with the frozen V39 train-plus-validation universe of 9,652 samples. A fixed graph joins exact decoded RGB matches, pHash/dHash candidate relations under the locked audit rule, and pairs from clusters that were manually reviewed as adjacent-or-near-identical. Connected components are assigned wholly to train or validation by a deterministic objective that targets the original validation count, minimizes moved samples, minimizes validation ground-truth-box difference, and then applies a stable lexical tie-break. No detection metrics, predictions, losses, or training results are used for assignment.

The accepted manifests contain 7,439 training images with 22,480 boxes and 2,213 validation images with 5,867 boxes. The unchanged 837-image guard partition is archival and non-test; it is excluded from model selection, performance reporting, and manuscript metrics. The post-assignment audit reports no train-validation sample overlap and no cross-boundary relation under the locked exact, perceptual, reviewed-adjacency, or extended-component audits. This supports a fixed-rule component-disjoint statement, not a universal claim that every possible scene dependence is absent.

\section{Experiments and Results}
\subsection{Frozen Training and Evaluation Protocol}
Both model groups use the same V40-v2 manifests, 50 epochs, $640\times640$ inputs, batch size 4, learning rate $10^{-4}$, AdamW, no mosaic/mixup/crop/flip/color-jitter augmentation, and fixed seeds 0 and 2. The detector score threshold is 0.001, the NMS threshold is 0.6, and at most 100 detections are retained per image. Precision, recall, and F1 use an operating threshold of 0.50. AP50 and AP75 are project-local single-class score-ranked AP values based on greedy one-to-one matching; they are not COCO AP50:95.

Within each run, the frozen trainer retains the checkpoint with the highest in-training validation AP50. The evidence is therefore validation-model-selection evidence, not an independent test. No V40 hyperparameter sweep is performed.

\subsection{Full-Input Comparison}
The pre-specified reliability-aware $p=0.15$ system has higher two-run mean precision, recall, F1, AP50, and AP75 than matched early fusion. The AP50, AP75, and F1 differences are $+0.0177$, $+0.0552$, and $+0.0188$, respectively. The reliability system has a smaller two-run AP50 range, but two seeds are not enough for a broad stability claim.

\begin{table*}[t]
\centering
\caption{V40-v2 full-input validation results. The reliability-aware $p=0.15$ setting was pre-specified before V40 results; no V40 dropout sweep or optimal-dropout claim is made. Range is the minimum--maximum range across two fixed seeds.}
\label{tab:core}
\scriptsize
\begin{tabular}{@{}llccccc@{}}
\toprule
Model group & Statistic & Precision & Recall & F1 & AP50 & AP75\\
\midrule
Matched early fusion & seed 0 & 0.9092 & 0.8955 & 0.9023 & 0.9455 & 0.8411\\
Matched early fusion & seed 2 & 0.9284 & 0.8485 & 0.8866 & 0.9361 & 0.8004\\
Matched early fusion & mean & 0.9188 & 0.8720 & 0.8945 & 0.9408 & 0.8208\\
Matched early fusion & range & 0.0192 & 0.0470 & 0.0157 & 0.0094 & 0.0407\\
\midrule
Reliability-aware $p=0.15$ & seed 0 & 0.9333 & 0.8972 & 0.9149 & 0.9584 & 0.8952\\
Reliability-aware $p=0.15$ & seed 2 & 0.9213 & 0.9022 & 0.9116 & 0.9588 & 0.8567\\
Reliability-aware $p=0.15$ & mean & \textbf{0.9273} & \textbf{0.8997} & \textbf{0.9133} & \textbf{0.9586} & \textbf{0.8760}\\
Reliability-aware $p=0.15$ & range & 0.0120 & 0.0049 & 0.0033 & 0.0004 & 0.0385\\
\bottomrule
\end{tabular}
\end{table*}

\begin{figure}[t]
\centering
\includegraphics[width=\columnwidth]{figures/fig3_core_results.pdf}
\caption{Two-run means for full-input V40-v2 validation. Error bars are population standard deviations over the two fixed seeds and are descriptive only.}
\label{fig:core}
\end{figure}

\subsection{Image-Level Bootstrap Uncertainty}
V40 validation images are resampled 2,000 times with a fixed bootstrap seed. Each resample computes the metric for each fixed checkpoint, averages the two checkpoints within a model group, and calculates the reliability-minus-early difference. Images with no ground-truth boxes remain in the resampling frame. All percentile intervals for AP50, AP75, and F1 are positive: AP50 $+0.0177$ [$+0.0153$, $+0.0203$], AP75 $+0.0553$ [$+0.0483$, $+0.0622$], and F1 $+0.0187$ [$+0.0150$, $+0.0226$]. These intervals are conditional on the fixed checkpoints and validation partition and are not used for model selection.

\subsection{Synthetic Channel Removal}
One channel group is deterministically set to zero at a time while model weights, thresholds, and postprocessing are fixed. Reliability-aware fusion retains higher AP50 under all three removals: RGB removed, 0.8985 versus 0.7379; thermal removed, 0.7030 versus 0.3648; event removed, 0.9582 versus 0.9339. This deterministic zero-channel intervention does not emulate measured physical sensor outages, temporal misalignment, or cross-device deployment failures.

\begin{figure}[t]
\centering
\includegraphics[width=\columnwidth]{figures/fig4_channel_removal.pdf}
\caption{AP50 under deterministic synthetic channel removal. Reliability-aware $p=0.15$ retains higher AP50 under all four input conditions.}
\label{fig:removal}
\end{figure}

\subsection{Efficiency}
The reliability-aware model adds 1,684 parameters and approximately 0.77 GFLOPs under the reported profiling convention. With batch size 1 and image size 640, its mean median raw-forward latency is 10.61 ms compared with 10.43 ms for early fusion, and its detector-inference latency is 23.30 ms compared with 22.89 ms. Its observed peak CUDA-memory range is also higher. Raw-forward timing covers the backbone call only; detector inference includes transform, backbone, head, NMS, and postprocessing, but excludes data loading and file I/O. Thus the reliability-aware system should not be described as faster.

\section{Discussion and Limitations}
The V40-v2 results support a narrow system-level conclusion: under the fixed V40-v2 component-disjoint validation protocol, the pre-specified reliability-aware $p=0.15$ system outperforms matched early fusion in the two fixed training runs and under deterministic removal of a single modality group. The comparison does not isolate the causal contribution of the gate alone because the reliability model is trained with modality dropout while early fusion is not. The absence of trained RGB-only, thermal-only, event-only, and static-global-weight controls also limits what can be concluded about individual modalities or the specific mechanism of input-conditioned weighting.

The split protocol is stronger than an exact-duplicate check but has a defined scope. It enforces component disjointness for the locked graph composed of exact RGB content, perceptual-hash candidates, and human-adjudicated adjacent-or-near-identical observations. It cannot demonstrate that all possible scene correlations are absent. Filename identifiers were not treated as verified temporal metadata. The guard partition is not a test set and is not used to make any performance claim.

The experimental evidence is validation-only. The V40 validation partition participates in within-run checkpoint retention and final reporting. A separately acquired and locked tri-modal test set would be required for independent-test or external-generalization claims. Synthetic channel removal is a controlled zeroing operation, not an estimate of real physical sensor reliability.

\section{Conclusion}
RA-RepDet is a lightweight reliability-aware RGB-thermal-event detector evaluated against matched early fusion on a V40-v2 expanded-adjacency component-disjoint validation partition. The pre-specified reliability-aware $p=0.15$ system improves two-run mean AP50, AP75, F1, precision, and recall, and its advantage remains visible under deterministic synthetic removal of RGB, thermal, or event inputs. The model carries modest parameter and compute overhead but higher peak memory and slightly higher measured latency. The contribution is therefore a validation-only, reproducible fusion result with a transparent partition-audit scope, rather than a claim of independent testing, universal leakage absence, or physical sensor-failure robustness.

\begin{thebibliography}{99}
\bibitem{zhao2019object} Zhao, Z.-Q., Zheng, P., Xu, S.-T., Wu, X.: Object detection with deep learning: A review. IEEE Trans. Neural Netw. Learn. Syst. 30(11), 3212--3232 (2019). https://doi.org/10.1109/TNNLS.2018.2876865
\bibitem{cheng2016survey} Cheng, G., Han, J.: A survey on object detection in optical remote sensing images. ISPRS J. Photogramm. Remote Sens. 117, 11--28 (2016). https://doi.org/10.1016/j.isprsjprs.2016.03.014
\bibitem{xia2022dota} Xia, G.-S. et al.: DOTA: A large-scale dataset for object detection in aerial images. IEEE Trans. Pattern Anal. Mach. Intell. 44(10), 5969--5988 (2022). https://doi.org/10.1109/TPAMI.2021.3074724
\bibitem{zhu2017deepRS} Zhu, X.X. et al.: Deep learning in remote sensing: A comprehensive review and list of resources. IEEE Geosci. Remote Sens. Mag. 5(4), 8--36 (2017). https://doi.org/10.1109/MGRS.2017.2762307
\bibitem{baltrusaitis2019multimodal} Baltru\v{s}aitis, T., Ahuja, C., Morency, L.-P.: Multimodal machine learning: A survey and taxonomy. IEEE Trans. Pattern Anal. Mach. Intell. 41(2), 423--443 (2019). https://doi.org/10.1109/TPAMI.2018.2798607
\bibitem{gallego2022event} Gallego, G. et al.: Event-based vision: A survey. IEEE Trans. Pattern Anal. Mach. Intell. 44(1), 154--180 (2022). https://doi.org/10.1109/TPAMI.2020.3008413
\bibitem{kapoor2023leakage} Kapoor, S., Narayanan, A.: Leakage and the reproducibility crisis in machine-learning-based science. Patterns 4(9), 100804 (2023). https://doi.org/10.1016/j.patter.2023.100804
\bibitem{cawley2010overfitting} Cawley, G.C., Talbot, N.L.C.: On over-fitting in model selection and subsequent selection bias in performance evaluation. J. Mach. Learn. Res. 11, 2079--2107 (2010)
\bibitem{wang2024repvit} Wang, A., Chen, H., Lin, Z., Han, J., Ding, G.: RepViT: Revisiting mobile CNN from ViT perspective. In: Proc. IEEE/CVF Conf. Comput. Vis. Pattern Recognit., pp. 15909--15920 (2024)
\bibitem{lin2017fpn} Lin, T.-Y. et al.: Feature pyramid networks for object detection. In: Proc. IEEE Conf. Comput. Vis. Pattern Recognit., pp. 2117--2125 (2017). https://doi.org/10.1109/CVPR.2017.106
\bibitem{tian2019fcos} Tian, Z., Shen, C., Chen, H., He, T.: FCOS: Fully convolutional one-stage object detection. In: Proc. IEEE/CVF Int. Conf. Comput. Vis., pp. 9627--9636 (2019). https://doi.org/10.1109/ICCV.2019.00972
\bibitem{ren2017faster} Ren, S., He, K., Girshick, R., Sun, J.: Faster R-CNN: Towards real-time object detection with region proposal networks. IEEE Trans. Pattern Anal. Mach. Intell. 39(6), 1137--1149 (2017). https://doi.org/10.1109/TPAMI.2016.2577031
\bibitem{lin2020focal} Lin, T.-Y., Goyal, P., Girshick, R., He, K., Doll\'{a}r, P.: Focal loss for dense object detection. IEEE Trans. Pattern Anal. Mach. Intell. 42(2), 318--327 (2020). https://doi.org/10.1109/TPAMI.2018.2858826
\bibitem{ramachandram2017deep} Ramachandram, D., Taylor, G.W.: Deep multimodal learning: A survey on recent advances and trends. IEEE Signal Process. Mag. 34(6), 96--108 (2017). https://doi.org/10.1109/MSP.2017.2738401
\bibitem{li2017pixel} Li, S., Kang, X., Fang, L., Hu, J., Yin, H.: Pixel-level image fusion: A survey of the state of the art. Inf. Fusion 33, 100--112 (2017). https://doi.org/10.1016/j.inffus.2016.05.004
\bibitem{ma2019surveyfusion} Ma, J., Ma, Y., Li, C.: Infrared and visible image fusion methods and applications: A survey. Inf. Fusion 45, 153--178 (2019). https://doi.org/10.1016/j.inffus.2018.02.004
\bibitem{li2019densefuse} Li, H., Wu, X.-J.: DenseFuse: A fusion approach to infrared and visible images. IEEE Trans. Image Process. 28(5), 2614--2623 (2019). https://doi.org/10.1109/TIP.2018.2887342
\bibitem{lichtsteiner2008sensor} Lichtsteiner, P., Posch, C., Delbr\"uck, T.: A 128$\times$128 120 dB 15 $\mu$s latency asynchronous temporal contrast vision sensor. IEEE J. Solid-State Circuits 43(2), 566--576 (2008). https://doi.org/10.1109/JSSC.2007.914337
\bibitem{gehrig2021dsec} Gehrig, M., Aarents, W., Gehrig, D., Scaramuzza, D.: DSEC: A stereo event camera dataset for driving scenarios. IEEE Robot. Autom. Lett. 6(3), 4947--4954 (2021). https://doi.org/10.1109/LRA.2021.3068942
\bibitem{neverova2016moddrop} Neverova, N., Wolf, C., Taylor, G.W., Nebout, F.: ModDrop: Adaptive multi-modal gesture recognition. IEEE Trans. Pattern Anal. Mach. Intell. 38(8), 1692--1706 (2016). https://doi.org/10.1109/TPAMI.2015.2461544
\end{thebibliography}


\end{document}
